\section{Steady-State Computation and Log-Linearization}
\label{sec:steady-state-computation}

This appendix provides the derivation of the steady-state values for capital, dividends, and debt, and the log-linearization of the flow of funds constraint using the first-order Taylor expansion.

\subsection{Steady-State Computation}
Notation marks:
\begin{itemize}
    \item lower case letter represents log deviation from the steady state
    \item upper case letter represents a level variable
    \item \(X^*\) represents the optimal path
\end{itemize}
The utility function is logarithmic:
\begin{equation}
U(D) = \ln(D)
\end{equation}

The flow of funds constraint is:
\begin{equation}
K^{t+1} + \left( \frac{R_f}{p} - \frac{(1-p)}{p} \frac{\mu Z K_t^\alpha}{B} \right) B_t + D_t = Z K^\alpha + B^{t+1}
\end{equation}

At the steady state, we have:
\begin{equation}
K + \left( \frac{R_f}{p} - \frac{(1-p)}{p} \frac{\mu Z K^\alpha}{B} \right) B + D = Z K^\alpha + B
\end{equation}

The Lagrangian for the firm's problem is:
\begin{equation}
\mathcal{L} = \sum_{t=0}^\infty \beta^t \left[ \ln(D_t) - \lambda_t \left( - Z K_t^\alpha - B_{t+1} + K_{t+1} + \left( \frac{R_f}{p} - \frac{(1-p)}{p} \frac{\mu Z K_t^\alpha}{B_t} \right) B_t + D_t \right) \right]
\end{equation}

the first-order conditions are:
\begin{itemize}
    \item For Dividends:
    \begin{equation}
        \frac{\partial \mathcal{L}}{\partial D_t} = \frac{1}{D_t} - \lambda_t = 0 \implies \lambda_t = \frac{1}{D_t}
    \end{equation}
    \item For Capital:
    \begin{align}
        &\frac{\partial \mathcal{L}}{\partial K_{t+1}} = -\lambda_t + \beta \lambda_{t+1} \alpha Z K_{t+1}^{\alpha-1} 
        \left(\frac{p+(1-p)\mu}{p}\right) = 0 \\
        & \implies \lambda_t = \beta \lambda_{t+1} \alpha Z K_{t+1}^{\alpha-1} \left(\frac{p+(1-p)\mu}{p}\right)
    \end{align}
    \item For Debt:
    \begin{equation}
        \frac{\partial \mathcal{L}}{\partial B_{t+1}} = \lambda_t - \beta \lambda_{t+1} \frac{R_f}{p}= 0 \implies \lambda_t = \beta \lambda_{t+1} \frac{R_f}{p}
    \end{equation}
\end{itemize}

Combining the first-order conditions, we get the Euler equations for capital and debt.

\paragraph{Euler Equation for Capital}
\begin{equation}
\frac{1}{D_t} = \beta \frac{1}{D_{t+1}} \alpha Z K_{t+1}^{\alpha-1} \left(\frac{p+(1-p)\mu}{p}\right) \label{eq:euler-capital}
\end{equation}

\paragraph{Euler Equation for Debt}
\begin{equation}
\frac{1}{D_t} = \beta \frac{1}{D_{t+1}} \frac{R_f}{p} \label{eq:euler-debt}
\end{equation}

At the steady state, \(K_t = K_{t+1} = K\), \(D_t = D_{t+1} = D\), and \(B_t = B_{t+1} = B\).

\subsubsection{Steady-State Euler Equation for Capital}
\begin{equation}
K = \left(\frac{\alpha \beta Z (p+(1-p)\mu)}{p}\right)^{\frac{1}{1-\alpha}}
\end{equation}

\subsubsection{Steady-State Euler Equation for Debt}
\begin{equation}
p = \frac{\beta}{R_f}
\end{equation}


\subsubsection{Steady-State Flow of Funds Constraint}
Using the flow of funds constraint at the steady state:
\begin{equation}
K + \left( \frac{R_f}{p} - \frac{(1-p)}{p} \frac{\mu Z K^\alpha}{B} \right) B + D = Z K^\alpha + B
\end{equation}

we get the steady state level of dividends:
\begin{equation}
D = Z K^\alpha \left(\frac{p + \mu (1-p)}{p}\right) - K - B \left(\frac{R_f - p}{p}\right)
\end{equation}

\subsection{Log-Linearization Using Taylor Expansion}

We denote deviations from the steady state using lower-case letters for the log variables and upper-case letters for the level variables.

\subsubsection{Flow of Funds Constraint}

The original flow of funds constraint is:
\begin{equation}
K_{t+1} + \left( \frac{R_f}{p} - \frac{(1-p)}{p} \frac{\mu F(K_t)}{B_t} \right) B_t + D_t = F(K_t) + B_{t+1}
\end{equation}
where $F(K_t) = Z K_t^\alpha$.

At the steady state, we have:
\begin{equation}
K + \left( \frac{R_f}{p} - \frac{(1-p)}{p} \frac{\mu Z K^\alpha}{B} \right) B + D = Z K^\alpha + B
\end{equation}

Using log deviations:
\begin{equation}
k_t = \ln\left(\frac{K_t}{K}\right), \quad d_t = \ln\left(\frac{D_t}{D}\right), \quad b_t = \ln\left(\frac{B_t}{B}\right), \quad k_{t+1} = \ln\left(\frac{K_{t+1}}{K}\right), \quad b_{t+1} = \ln\left(\frac{B_{t+1}}{B}\right)
\end{equation}

\subsubsection{Linearizing the Flow of Funds Constraint}


Using a first-order Taylor expansion around the steady state (\(K, B, D\)), we get:
\begin{equation}
    K_{t+1} = f(K) + (1-R)B - D + \mathbb{J}^{T}\begin{bmatrix}
        K_t-K \\
        B_t-B \\
        B_{t+1}-B \\
        D_t-D
\end{bmatrix}
\end{equation}

where:
\begin{equation}
    \mathbb{J} = \begin{bmatrix}
        \frac{\partial K_{t+1}}{\partial K_t} \\
        \frac{\partial K_{t+1}}{\partial B_t}\\
        \frac{\partial K_{t+1}}{\partial B_{t+1}}\\
        \frac{\partial K_{t+1}}{\partial D_t}
    \end{bmatrix}
\end{equation}
exploiting the Euler equations for capital and debt \ref{eq:euler-capital} and \ref{eq:euler-debt} we get:
\begin{equation}
    (K_{t+1}-K) - (B_{t+1}-B) = \frac{1}{\beta} \left[\left(K_t - K\right) - \left(B_t - B\right)\right] - \left(D_t - D\right)
\end{equation}
rewriting using lowercase letters to express the log deviation from the steady state:
\begin{equation}
    k_{t+1} - b{t+1} = \frac{1}{\beta} \left(k_t - b_t\right) - d_t
\end{equation}
we can rewrite expressing the equity as (\(n=k-b\)):
\begin{equation}
    n_{t+1} = \frac{1}{\beta} n_t - d_t
\end{equation}
This completes the log-linearization of the flow of funds constraint around the steady state.
